\documentclass[12pt, a4paper]{article}

% 1. Cấu hình Tiếng Việt và Font Times New Roman 13pt chuẩn
\usepackage[utf8]{inputenc}
\usepackage[T5]{fontenc}
\usepackage[vietnamese]{babel}
\usepackage{newtxtext,newtxmath}

\makeatletter
\renewcommand{\normalsize}{\@setfontsize\normalsize{13pt}{16pt}}
\makeatother
\normalsize

% 2. Gói Toán học & Ký hiệu bổ trợ
\usepackage{amsmath, amssymb, amsfonts, mathrsfs, amsthm}
\usepackage{graphicx}
\usepackage{fontawesome}
\usepackage{booktabs, tabularx, array, multirow}
\usepackage{enumitem}

% 3. Đồ họa TikZ, Bảng biến thiên & Biểu đồ
\usepackage{tikz}
\usetikzlibrary{calc, intersections, angles, quotes, patterns, arrows.meta, 3d}
\usepackage{pgfplots}
\pgfplotsset{compat=1.18}
\usepackage{tkz-euclide}
\usepackage{tkz-tab}

\renewcommand{\vec}[1]{\overrightarrow{#1}}

% 4. Hộp sư phạm (Tcolorbox)
\usepackage[many]{tcolorbox}
\tcbuselibrary{skins, breakable}

\newtcolorbox{pedabox}[1][]{
    enhanced,
    colback=blue!5!white,
    colframe=blue!75!black,
    fonttitle=\bfseries\sffamily,
    title=#1,
    boxrule=0.8pt,
    arc=2mm,
    breakable,
    left=3mm, right=3mm, top=2mm, bottom=2mm
}

\newtcolorbox{scaffoldbox}[1][]{
    enhanced,
    colback=orange!5!white,
    colframe=orange!85!black,
    fonttitle=\bfseries\sffamily,
    title=#1,
    boxrule=0.8pt,
    arc=2mm,
    breakable,
    left=3mm, right=3mm, top=2mm, bottom=2mm
}

\newtcolorbox{aibox}[1][]{
    enhanced,
    colback=green!5!white,
    colframe=teal!80!black,
    fonttitle=\bfseries\sffamily,
    title=#1,
    boxrule=0.8pt,
    arc=2mm,
    breakable,
    left=3mm, right=3mm, top=2mm, bottom=2mm
}

% 5. Căn lề chuẩn văn bản giáo án
\usepackage{geometry}
\geometry{
    a4paper,
    left=25mm,
    right=15mm,
    top=20mm,
    bottom=20mm
}

% 6. Header / Footer chuẩn hóa (BẮT BUỘC CHÍNH XÁC NỘI DUNG NÀY)
\usepackage{fancyhdr}
\usepackage{lastpage}
\pagestyle{fancy}
\fancyhf{}

\fancyhead[L]{\small \textit{Kế hoạch bài dạy Toán 11}}
\fancyhead[R]{\textbf{Trang \thepage/\pageref{LastPage}}}
\fancyfoot[L]{\small \textit{Tổ Toán - Tin}}
\fancyfoot[R]{\small \textit{Trường THCS\&THPT Hồng Vân}}
\renewcommand{\headrulewidth}{0.4pt}
\renewcommand{\footrulewidth}{0.4pt}

% 7. Giãn dòng & Giãn đoạn theo chuẩn Nghị định 30/2020/NĐ-CP
\usepackage{setspace}
\setstretch{1.15}
\setlength{\parindent}{1.0cm}
\setlength{\parskip}{5pt}

\begin{document}

\begin{center}
    {\Large \textbf{KẾ HOẠCH BÀI DẠY MÔN TOÁN LỚP 11}}\\[0.4em]
    {\LARGE \textbf{BÀI TẬP CUỐI CHƯƠNG III}}\\[0.3em]
    {\large \textbf{(CÁC SỐ ĐẶC TRƯNG ĐO XU THẾ TRUNG TÂM}}\\[0.2em]
    {\large \textbf{CỦA MẪU SỐ LIỆU GHÉP NHÓM)}}\\[0.5em]
    \textbf{Môn học:} Toán 11 | \textbf{Thời lượng:} 01 tiết (Tiết 21 theo PPCT)
\end{center}

\vspace{0.5em}
\hrule
\vspace{1em}

\section*{I. MỤC TIÊU DẠY HỌC}

\subsection*{1. Kiến thức, Năng lực}

\textbf{a) Năng lực Toán học đặc thù (nguyên văn chuẩn CT GDPT 2018):}
\begin{itemize}[leftmargin=1.5em]
    \item \textbf{Năng lực tư duy và lập luận toán học:}
    \begin{itemize}
        \item Hệ thống hóa và củng cố toàn diện các số đặc trưng đo xu thế trung tâm của mẫu số liệu ghép nhóm: số trung bình cộng ($\overline{x}$), trung vị ($M_e$), các tứ phân vị ($Q_1, Q_2, Q_3$), khoảng tứ phân vị ($\Delta_Q = Q_3 - Q_1$), mốt ($M_o$).
        \item So sánh, phân tích ưu điểm, nhược điểm và phạm vi áp dụng tối ưu của từng số đặc trưng:
        \begin{itemize}
            \item Số trung bình $\overline{x}$ sử dụng toàn bộ thông tin của mẫu nhưng rất nhạy cảm với các giá trị bất thường (outliers).
            \item Trung vị $M_e$ và khoảng tứ phân vị $\Delta_Q$ không bị ảnh hưởng bởi các giá trị bất thường, thích hợp cho phân phối lệch.
            \item Mốt $M_o$ phản ánh giá trị có mật độ quan sát dày đặc nhất, đặc biệt hữu ích trong kinh doanh và quản lý.
        \end{itemize}
        \item Rút ra kết luận có ý nghĩa thực tiễn từ các số đặc trưng đo xu thế trung tâm khi so sánh hai mẫu số liệu ghép nhóm cùng loại.
    \end{itemize}
    \item \textbf{Năng lực giải quyết vấn đề toán học:}
    \begin{itemize}
        \item Thực hiện thành thạo quy trình tính số trung bình cộng $\overline{x}$ qua bảng 4 cột (Nhóm $\to$ Tần số $\to$ Giá trị đại diện $\to$ Tích số).
        \item Xác định chính xác nhóm chứa mốt và vận dụng công thức nội suy tính mốt $M_o$.
        \item Lập bảng tần số tích lũy $cf_i$, xác định đúng nhóm chứa trung vị, nhóm chứa tứ phân vị $Q_1, Q_3$ và tính chính xác các giá trị $M_e, Q_1, Q_3, \Delta_Q$.
        \item Giải quyết bài toán thực tế so sánh sự đồng đều và xu thế trung tâm giữa hai tập dữ liệu khảo sát.
    \end{itemize}
    \item \textbf{Năng lực mô hình hóa toán học:}
    \begin{itemize}
        \item Vận dụng các số đặc trưng đo xu thế trung tâm vào các bài toán phân tích thực tế: so sánh kết quả học tập của học sinh, năng suất lao động của các dây chuyền sản xuất, chỉ số thể lực học đường.
    \end{itemize}
    \item \textbf{Năng lực giao tiếp toán học:}
    \begin{itemize}
        \item Sử dụng chính xác các thuật ngữ toán học; trình bày lời giải rõ ràng, mạch lạc; phân tích và tranh luận logic trong hoạt động nhóm.
    \end{itemize}
    \item \textbf{Năng lực sử dụng công cụ, phương tiện học toán:}
    \begin{itemize}
        \item Khai thác tối ưu máy tính cầm tay (MTCT Casio fx-580VN X) tính toán phân số và các biểu thức nội suy phức tạp; sử dụng phần mềm bảng tính Excel đối chiếu kết quả.
    \end{itemize}
\end{itemize}

\textbf{b) Năng lực số (theo Thông tư 02/2025/TT-BGDĐT):}
\begin{itemize}[leftmargin=1.5em]
    \item \textbf{Mã 4.1NC1a (Bảo vệ dữ liệu cá nhân và quyền riêng tư trong môi trường số):} Học sinh nhận thức và thực hiện nghiêm túc việc bảo mật dữ liệu cá nhân khi làm việc với các bảng số liệu khảo sát học đường; thực hiện ẩn danh hóa thông tin (loại bỏ họ tên, số định danh cá nhân), không tùy tiện chia sẻ dữ liệu khảo sát nhạy cảm lên mạng xã hội, tuân thủ nguyên tắc đạo đức và an toàn thông tin số.
\end{itemize}

\textbf{c) Năng lực AI (theo Quyết định 2422/QĐ-BGDĐT):}
\begin{itemize}[leftmargin=1.5em]
    \item \textbf{Mã 11.A3.1 (Quyền của người dùng dữ liệu và đạo đức trong hệ thống AI):} Học sinh trình bày được quyền của người dùng dữ liệu (quyền được biết, quyền đồng ý, quyền yêu cầu chỉnh sửa hoặc xóa bỏ dữ liệu) khi các nền tảng và hệ thống trí tuệ nhân tạo (AI) thu thập thông tin hành vi người dùng; nhận diện nguy cơ định kiến thuật toán (Algorithmic Bias) khi tập dữ liệu huấn luyện thống kê bị thiên lệch hoặc thiếu tính đại diện.
\end{itemize}

\textbf{d) Năng lực chung \& Phẩm chất:}
\begin{itemize}[leftmargin=1.5em]
    \item \textbf{Năng lực chung:} Năng lực tự chủ và tự học (tự hệ thống kiến thức toàn chương, tự rà soát kĩ năng cá nhân); năng lực giao tiếp và hợp tác (thảo luận nhóm so sánh dữ liệu).
    \item \textbf{Phẩm chất:} Trung thực, khách quan trong xử lý số liệu; cẩn thận, chính xác trong tính toán; có trách nhiệm với cộng đồng trong quản lý và bảo vệ dữ liệu số.
\end{itemize}

\section*{II. THIẾT BỊ DẠY HỌC VÀ HỌC LIỆU}
\begin{itemize}[leftmargin=1.5em]
    \item \textbf{Giáo viên:} Kế hoạch bài dạy chuẩn CV 5512; bài giảng PowerPoint tóm tắt sơ đồ tư duy Chương III; Phiếu học tập số 1 (Bảng tổng hợp công thức), Phiếu học tập số 2 (Bài tập luyện tập phân hóa); máy tính, máy chiếu.
    \item \textbf{Học sinh:} Vở ghi bài, SGK Toán 11; máy tính cầm tay Casio fx-580VN X; thiết bị học tập.
\end{itemize}

\section*{III. TIẾN TRÌNH DẠY HỌC (01 TIẾT --- 45 PHÚT)}

\begin{center}
    \framebox[1.0\textwidth]{\parbox{0.96\textwidth}{
        \textbf{CẤU TRÚC TIẾN TRÌNH TIẾT 21 (BÀI TẬP CUỐI CHƯƠNG III):}\\[0.3em]
        $\bullet$ \textbf{Hoạt động 1 (05 phút):} Khởi động --- Trò chơi "Ghép đôi hoàn hảo" đối chiếu tên số đặc trưng và ý nghĩa thực tế.\\[0.3em]
        $\bullet$ \textbf{Hoạt động 2 (12 phút):} Hệ thống hóa kiến thức trọng tâm Chương III --- Giàn giáo sư phạm HS TB - Yếu --- Tích hợp Năng lực số (Mã 4.1NC1a) \& Năng lực AI (Mã 11.A3.1).\\[0.3em]
        $\bullet$ \textbf{Hoạt động 3 (20 phút):} Luyện tập giải toán phân hóa (3 dạng bài tập cốt lõi).\\[0.3em]
        $\bullet$ \textbf{Hoạt động 4 (08 phút):} Vận dụng thực tiễn --- Phân tích dữ liệu khảo sát và an toàn dữ liệu số.
    }}
\end{center}

\vspace{0.8em}
\hrule
\vspace{0.8em}

\subsubsection*{1. Hoạt động 1: Mở đầu / Khởi động (05 phút)}
\textbf{Chủ đề:} \textit{Trò chơi "Ghép đôi hoàn hảo": Số đặc trưng nào cho tình huống nào?}
\begin{itemize}
    \item \textbf{a) Mục tiêu:} Tái hiện nhanh vai trò, ý nghĩa của các số đặc trưng đo xu thế trung tâm; kích thích tư duy so sánh và lựa chọn số đo phù hợp với bối cảnh thực tiễn.
    \item \textbf{b) Nội dung:} Giáo viên đưa ra 4 số đặc trưng ở cột A và 4 tình huống thực tiễn ở cột B, yêu cầu học sinh ghép nối:
    \begin{itemize}
        \item Cột A: (1) Số trung bình $\overline{x}$; (2) Trung vị $M_e$; (3) Mốt $M_o$; (4) Khoảng tứ phân vị $\Delta_Q$.
        \item Cột B:
        \begin{itemize}
            \item (P) Chủ cửa hàng quần áo muốn nhập thêm cỡ áo bán chạy nhất.
            \item (Q) Nhà trường muốn tính điểm thi bình quân của toàn bộ học sinh khối 11.
            \item (R) Chuyên gia kinh tế đánh giá mức thu nhập tiêu biểu của người dân khi xã hội có sự chênh lệch giàu nghèo cực lớn (xuất hiện tỉ phú).
            \item (S) Nhà phân tích tài chính muốn đo độ phân tán của $50\%$ mức chi tiêu ở phân khúc tầm trung.
        \end{itemize}
    \end{itemize}
    \item \textbf{c) Sản phẩm:} Kết quả ghép nối của học sinh:
    \begin{itemize}
        \item (1) --- (Q): Số trung bình dùng để tính giá trị bình quân của toàn thể.
        \item (2) --- (R): Trung vị đại diện trung thực cho thu nhập khi có giá trị ngoại lai cực lớn.
        \item (3) --- (P): Mốt xác định giá trị xuất hiện phổ biến nhất (cỡ áo bán chạy nhất).
        \item (4) --- (S): Khoảng tứ phân vị đo độ phân tán của $50\%$ số liệu chính giữa.
    \end{itemize}
    \item \textbf{d) Tổ chức thực hiện:} GV giao nhiệm vụ $\to$ HS suy nghĩ giơ bảng trả lời $\to$ GV nhận xét, chốt lại ý nghĩa thực tiễn và dẫn dắt vào bài học ôn tập cuối chương.
\end{itemize}

\subsubsection*{2. Hoạt động 2: Hệ thống hóa kiến thức trọng tâm (12 phút)}

\paragraph{Mục 1: Bảng tổng hợp công thức các số đặc trưng đo xu thế trung tâm}

\begin{center}
\begin{tabularx}{\linewidth}{|p{3.2cm}|X|p{3.5cm}|}
\hline
\textbf{Số đặc trưng} & \textbf{Công thức toán học} & \textbf{Điều kiện / Lưu ý} \\
\hline
\textbf{1. Số trung bình} & $\overline{x} = \dfrac{1}{n} \sum_{i=1}^k m_i c_i$ & $c_i = \dfrac{a_i + a_{i+1}}{2}$; $n = \sum m_i$. \\
\hline
\textbf{2. Mốt ($M_o$)} & $M_o = a_j + \dfrac{m_j - m_{j-1}}{(m_j - m_{j-1}) + (m_j - m_{j+1})} \cdot h$ & Nhóm $j$ có $m_j = \max$; $h = a_{j+1} - a_j$. \\
\hline
\textbf{3. Trung vị ($M_e = Q_2$)} & $M_e = a_p + \dfrac{\dfrac{n}{2} - cf_{p-1}}{m_p} \cdot h$ & Nhóm $p$ đầu tiên có $cf_p \ge \dfrac{n}{2}$. \\
\hline
\textbf{4. Tứ phân vị thứ nhất ($Q_1$)} & $Q_1 = a_p + \dfrac{\dfrac{n}{4} - cf_{p-1}}{m_p} \cdot h$ & Nhóm $p$ đầu tiên có $cf_p \ge \dfrac{n}{4}$. \\
\hline
\textbf{5. Tứ phân vị thứ ba ($Q_3$)} & $Q_3 = a_r + \dfrac{\dfrac{3n}{4} - cf_{r-1}}{m_r} \cdot h$ & Nhóm $r$ đầu tiên có $cf_r \ge \dfrac{3n}{4}$. \\
\hline
\end{tabularx}
\end{center}

\begin{scaffoldbox}[Giàn giáo sư phạm cho HS TB - Yếu: Mẹo tính trung vị và mốt không sai sót]
    \textbf{1. Tìm nhóm chứa trung vị / tứ phân vị:}\\
    $\bullet$ Tính ngay mốc vị trí: Trung vị là $\dfrac{n}{2}$, $Q_1$ là $\dfrac{n}{4}$, $Q_3$ là $\dfrac{3n}{4}$.\\
    $\bullet$ Luôn lập cột tần số tích lũy $cf_i$. Dò từ trên xuống dưới, nhóm nào có $cf$ \textbf{vượt qua hoặc chạm mốc đầu tiên} thì đó chính là nhóm cần tìm.\\
    $\bullet$ \textbf{Lỗi thường gặp:} Lấy nhầm $cf_p$ thay vì $cf_{p-1}$ (phải lấy tần số tích lũy của nhóm \textbf{ĐỨNG TRƯỚC}) và chia nhầm cho $n$ thay vì chia cho $m_p$ (tần số của chính nhóm đó).\\
    \textbf{2. Tìm nhóm chứa mốt:}\\
    $\bullet$ Nhìn cột tần số $m_i$, nhóm nào có số lượng lớn nhất thì đó là nhóm chứa mốt.
\end{scaffoldbox}

\paragraph{Mục 2: Tích hợp Năng lực số (Mã 4.1NC1a) \& Năng lực AI (Mã 11.A3.1)}

\begin{pedabox}[Năng lực số (Mã 4.1NC1a): Bảo vệ dữ liệu cá nhân trong khảo sát số]
    $\bullet$ \textbf{Nguyên tắc bảo vệ quyền riêng tư:}\\
    Khi tiến hành thu thập dữ liệu học sinh (chiều cao, cân nặng, điểm thi, thói quen sinh hoạt) để xử lý số liệu ghép nhóm, học sinh phải tuân thủ nguyên tắc ẩn danh hóa: Tuyệt đối không lưu trữ hay công khai cột Họ tên, Số điện thoại, Địa chỉ nhà ở đi kèm số liệu.\\
    $\bullet$ Nhận thức rõ hành vi chia sẻ bảng điểm cá nhân của người khác lên không gian mạng khi chưa được cho phép là vi phạm quyền riêng tư số.
\end{pedabox}

\begin{aibox}[Tích hợp Năng lực AI (QĐ 2422/QĐ-BGDĐT --- Mã 11.A3.1): Đạo đức dữ liệu và định kiến AI]
    $\bullet$ \textbf{Quyền của người dùng dữ liệu:}\\
    Trong kỉ nguyên trí tuệ nhân tạo, các ứng dụng AI liên tục thu thập dữ liệu hành vi người dùng để tính toán các số đặc trưng thống kê và tối ưu thuật toán. Người dùng có quyền được thông báo minh bạch, quyền cho phép hoặc từ chối thu thập dữ liệu và quyền yêu cầu xóa vĩnh viễn dữ liệu cá nhân.\\
    $\bullet$ \textbf{Định kiến thuật toán (Algorithmic Bias):} Nếu dữ liệu thống kê đầu vào bị thiếu tính đại diện (ví dụ mẫu khảo sát chỉ thu thập ở khu vực thành thị giàu có), mô hình AI huấn luyện sẽ đưa ra kết luận sai lệch, phân biệt đối xử với các nhóm đối tượng khác.
\end{aibox}

\subsubsection*{3. Hoạt động 3: Luyện tập giải toán phân hóa (20 phút)}

\paragraph{Dạng 1: Tính số trung bình cộng và mốt (06 phút)}
\begin{itemize}
    \item \textbf{Bài toán 1:} Bảng số liệu sau ghi lại cự li ném bóng (đơn vị: mét) của $40$ học sinh:
    \begin{center}
    \begin{tabular}{|c|c|c|c|c|c|}
    \hline
    \textbf{Cự li ném (m)} & $[8; 10)$ & $[10; 12)$ & $[12; 14)$ & $[14; 16)$ & $[16; 18]$ \\
    \hline
    \textbf{Tần số ($m_i$)} & $4$ & $10$ & $16$ & $7$ & $3$ \\
    \hline
    \end{tabular}
    \end{center}
    Tính số trung bình cộng $\overline{x}$ và mốt $M_o$ của mẫu số liệu ghép nhóm trên.
    \item \textit{Lời giải chi tiết:}\\
    - Lập bảng tính toán 4 cột:
    \begin{center}
    \begin{tabular}{|c|c|c|c|}
    \hline
    \textbf{Nhóm cự li} & \textbf{Tần số ($m_i$)} & \textbf{Giá trị đại diện ($c_i$)} & $m_i \cdot c_i$ \\
    \hline
    $[8; 10)$ & $4$ & $9$ & $36$ \\
    $[10; 12)$ & $10$ & $11$ & $110$ \\
    $[12; 14)$ & $16$ & $13$ & $208$ \\
    $[14; 16)$ & $7$ & $15$ & $105$ \\
    $[16; 18]$ & $3$ & $17$ & $51$ \\
    \hline
    \textbf{Tổng} & $n = 40$ & --- & $\sum = 510$ \\
    \hline
    \end{tabular}
    \end{center}
    Số trung bình cộng: $\overline{x} = \dfrac{510}{40} = 12{,}75\text{ (m)}$.\\
    - Tính mốt $M_o$:\\
    Nhóm có tần số lớn nhất là $[12; 14)$ với $m_3 = 16 \implies j = 3$.\\
    Các thông số: $a_3 = 12$, $h = 2$, $m_3 = 16$, $m_2 = 10$, $m_4 = 7$.\\
    $$M_o = 12 + \dfrac{16 - 10}{(16 - 10) + (16 - 7)} \cdot 2 = 12 + \dfrac{6}{6 + 9} \cdot 2 = 12 + \dfrac{12}{15} = 12{,}8\text{ (m)}$$
\end{itemize}

\paragraph{Dạng 2: Tính trung vị, tứ phân vị và khoảng tứ phân vị (07 phút)}
\begin{itemize}
    \item \textbf{Bài toán 2:} Sử dụng số liệu ở Bài toán 1, hãy tính trung vị $M_e$, các tứ phân vị $Q_1, Q_3$ và khoảng tứ phân vị $\Delta_Q$.
    \item \textit{Lời giải chi tiết:}\\
    Bảng tần số tích lũy: $cf_1 = 4$; $cf_2 = 14$; $cf_3 = 30$; $cf_4 = 37$; $cf_5 = 40$.\\
    $\bullet$ \textbf{1. Tính trung vị $M_e$:}
    $\dfrac{n}{2} = 20$. Nhóm đầu tiên có $cf \ge 20$ là nhóm $[12; 14)$ (vì $cf_3 = 30 \ge 20$).\\
    Ta có: $a_3 = 12$, $h = 2$, $m_3 = 16$, $cf_2 = 14$.\\
    $$M_e = 12 + \dfrac{20 - 14}{16} \cdot 2 = 12 + \dfrac{6 \cdot 2}{16} = 12 + 0{,}75 = 12{,}75\text{ (m)}$$
    $\bullet$ \textbf{2. Tính tứ phân vị thứ nhất $Q_1$:}
    $\dfrac{n}{4} = 10$. Nhóm đầu tiên có $cf \ge 10$ là nhóm $[10; 12)$ (vì $cf_2 = 14 \ge 10$).\\
    Ta có: $a_2 = 10$, $h = 2$, $m_2 = 10$, $cf_1 = 4$.\\
    $$Q_1 = 10 + \dfrac{10 - 4}{10} \cdot 2 = 10 + \dfrac{6 \cdot 2}{10} = 10 + 1{,}2 = 11{,}2\text{ (m)}$$
    $\bullet$ \textbf{3. Tính tứ phân vị thứ ba $Q_3$:}
    $\dfrac{3n}{4} = 30$. Nhóm đầu tiên có $cf \ge 30$ là nhóm $[12; 14)$ (vì $cf_3 = 30$).\\
    Ta có: $a_3 = 12$, $h = 2$, $m_3 = 16$, $cf_2 = 14$.\\
    $$Q_3 = 12 + \dfrac{30 - 14}{16} \cdot 2 = 12 + \dfrac{16}{16} \cdot 2 = 14\text{ (m)}$$
    $\bullet$ \textbf{4. Khoảng tứ phân vị:}
    $$\Delta_Q = Q_3 - Q_1 = 14 - 11{,}2 = 2{,}8\text{ (m)}$$
\end{itemize}

\paragraph{Dạng 3: Bài toán tổng hợp so sánh hai mẫu số liệu ghép nhóm (07 phút)}
\begin{itemize}
    \item \textbf{Bài toán 3:} Kết quả kiểm tra môn Toán của hai lớp 11A và 11B (mỗi lớp gồm $40$ học sinh) được cho bởi bảng ghép nhóm:
    \begin{center}
    \begin{tabular}{|c|c|c|c|c|c|}
    \hline
    \textbf{Khung điểm} & $[0; 4)$ & $[4; 6)$ & $[6; 8)$ & $[8; 10]$ & \textbf{Tổng} \\
    \hline
    \textbf{Lớp 11A ($m_i$)} & $2$ & $10$ & $20$ & $8$ & $n_A = 40$ \\
    \hline
    \textbf{Lớp 11B ($m_i$)} & $6$ & $6$ & $18$ & $10$ & $n_B = 40$ \\
    \hline
    \end{tabular}
    \end{center}
    Hãy tính số trung bình $\overline{x}$ và trung vị $M_e$ của mỗi lớp. Nhận xét về năng lực học tập môn Toán của hai lớp.
    \item \textit{Lời giải chi tiết:}\\
    - Giá trị đại diện các nhóm: $c_1 = 2$; $c_2 = 5$; $c_3 = 7$; $c_4 = 9$.\\
    $\bullet$ \textbf{Lớp 11A:}\\
    $\overline{x}_A = \dfrac{2 \cdot 2 + 10 \cdot 5 + 20 \cdot 7 + 8 \cdot 9}{40} = \dfrac{4 + 50 + 140 + 72}{40} = \dfrac{266}{40} = 6{,}65$.\\
    Tần số tích lũy lớp 11A: $2, 12, 32, 40$. Vị trí $n/2 = 20 \implies$ Nhóm chứa trung vị là $[6; 8)$.\\
    $$M_{e(A)} = 6 + \dfrac{20 - 12}{20} \cdot 2 = 6 + \dfrac{8}{20} \cdot 2 = 6 + 0{,}8 = 6{,}8$$
    $\bullet$ \textbf{Lớp 11B:}\\
    $\overline{x}_B = \dfrac{6 \cdot 2 + 6 \cdot 5 + 18 \cdot 7 + 10 \cdot 9}{40} = \dfrac{12 + 30 + 126 + 90}{40} = \dfrac{258}{40} = 6{,}45$.\\
    Tần số tích lũy lớp 11B: $6, 12, 30, 40$. Vị trí $n/2 = 20 \implies$ Nhóm chứa trung vị là $[6; 8)$.\\
    $$M_{e(B)} = 6 + \dfrac{20 - 12}{18} \cdot 2 = 6 + \dfrac{8}{18} \cdot 2 = 6 + \dfrac{8}{9} \approx 6{,}89$$
    - \textbf{Nhận xét so sánh:}
    Lớp 11A có điểm trung bình cao hơn lớp 11B ($6{,}65 > 6{,}45$) do lớp 11B có nhiều học sinh bị điểm yếu ($6$ em khung $[0; 4)$ so với $2$ em của 11A). Tuy nhiên, trung vị của 11B lại cao hơn ($6{,}89 > 6{,}80$) và số học sinh đạt điểm giỏi của lớp 11B cũng nhiều hơn ($10$ em so với $8$ em). Như vậy, lớp 11A học đều hơn, trong khi lớp 11B có sự phân hóa mạnh mẽ hơn giữa nhóm học sinh xuất sắc và nhóm học sinh yếu.
\end{itemize}

\subsubsection*{4. Hoạt động 4: Vận dụng thực tiễn (08 phút)}
\textbf{Chủ đề:} \textit{Khảo sát thời gian sử dụng điện thoại thông minh và đạo đức dữ liệu}
\begin{itemize}
    \item \textbf{a) Mục tiêu:} Vận dụng các số đặc trưng đánh giá thói quen sử dụng công nghệ số; nâng cao nhận thức bảo vệ dữ liệu cá nhân khi tham gia khảo sát trực tuyến.
    \item \textbf{b) Nội dung:} Nhóm học sinh thực hiện khảo sát thời gian sử dụng mạng xã hội hàng ngày của $60$ bạn học sinh. Khi lập bảng số liệu và báo cáo kết quả:
    \begin{itemize}
        \item 1. Cần tính các số đặc trưng nào để phản ánh đúng thực trạng mà không bị kéo lệch bởi những bạn sử dụng điện thoại cá biệt (trên $8$ giờ/ngày)?
        \item 2. Làm thế nào để đảm bảo tính ẩn danh và bảo vệ quyền riêng tư cho các bạn học sinh tham gia khảo sát?
    \end{itemize}
    \item \textbf{c) Sản phẩm:}\\
    - 1. Nên sử dụng \textbf{Trung vị} ($M_e$) và \textbf{Khoảng tứ phân vị} ($\Delta_Q$) vì các đại lượng này không bị ảnh hưởng bởi một vài bạn "nghiện" điện thoại có số giờ bất thường (giá trị ngoại lai).\\
    - 2. Để bảo vệ dữ liệu: Biểu mẫu khảo sát không được yêu cầu nhập họ tên hoặc địa chỉ email cá nhân; dữ liệu thô chỉ được lưu trữ nội bộ và tiêu hủy sau khi hoàn thành đề tài; báo cáo chỉ trình bày dưới dạng bảng tần số ghép nhóm tổng thể.
    \item \textbf{d) Tổ chức thực hiện:} GV nêu tình huống $\to$ Các nhóm thảo luận, cử đại diện trình bày $\to$ GV nhận xét và kết luận bài học.
\end{itemize}

\newpage

\section*{IV. PHỤ LỤC HỌC LIỆU BỔ TRỢ}

\subsection*{Phụ lục 1: Phiếu học tập số 1 --- Bảng cẩm nang tra cứu công thức Chương III}
\noindent\textbf{Họ và tên:} .................................................................... \textbf{Lớp:} 11A....... \textbf{Nhóm:} ..........

\begin{pedabox}[Nhiệm vụ: Điền thông tin hoàn thiện cẩm nang công thức]
\begin{center}
\begin{tabularx}{0.95\linewidth}{|l|X|l|}
\hline
\textbf{Tên đại lượng} & \textbf{Công thức tính} & \textbf{Điều kiện nhận diện nhóm} \\
\hline
Số trung bình $\overline{x}$ & $\overline{x} = $ ......................................................... & $c_i = $ ................................... \\
\hline
Mốt $M_o$ & $M_o = a_j + \dfrac{\dots\dots\dots\dots\dots}{\dots\dots\dots\dots\dots} \cdot h$ & $m_j = $ .................................. \\
\hline
Trung vị $M_e$ & $M_e = a_p + \dfrac{\dots\dots - cf_{p-1}}{m_p} \cdot h$ & $cf_p \ge $ .............................. \\
\hline
Tứ phân vị $Q_1$ & $Q_1 = a_p + \dfrac{\dots\dots - cf_{p-1}}{m_p} \cdot h$ & $cf_p \ge $ .............................. \\
\hline
Tứ phân vị $Q_3$ & $Q_3 = a_r + \dfrac{\dots\dots - cf_{r-1}}{m_r} \cdot h$ & $cf_r \ge $ .............................. \\
\hline
Khoảng tứ phân vị & $\Delta_Q = $ ..................................................... & Ý nghĩa: ................................. \\
\hline
\end{tabularx}
\end{center}
\end{pedabox}

\subsection*{Phụ lục 2: Phiếu học tập số 2 --- Bài tập tự luyện và câu hỏi tình huống AI}
\noindent\textbf{Họ và tên:} .................................................................... \textbf{Lớp:} 11A....... \textbf{Nhóm:} ..........

\begin{pedabox}[Nhiệm vụ: Luyện tập cá nhân]
$\bullet$ \textbf{Bài 1 (Cơ bản):} Cho mẫu số liệu ghép nhóm có $n = 50$, nhóm $[20; 30)$ có tần số $m = 15$, tần số tích lũy của nhóm trước là $cf = 18$. Tính trung vị $M_e$.\\
$\bullet$ \textbf{Bài 2 (Tình huống AI):} Tại sao khi các ứng dụng trí tuệ nhân tạo (AI) đề xuất thời gian giao hàng dự kiến cho khách hàng, thuật toán thường căn cứ vào trung vị $M_e$ thay vì số trung bình $\overline{x}$ của các chuyến xe trước đó?\\
\textit{Gợi ý trả lời:} ..................................................................................................................................
\end{pedabox}

\subsection*{Phụ lục 3: Bảng Rubric đánh giá năng lực tích hợp trong Bài tập cuối chương III}

\begin{center}
\begin{tabularx}{\linewidth}{|p{3.2cm}|X|X|X|X|}
\hline
\textbf{Tiêu chí} & \textbf{Mức 1 (Chưa đạt)} & \textbf{Mức 2 (Đạt)} & \textbf{Mức 3 (Khá)} & \textbf{Mức 4 (Tốt)} \\
\hline
\textbf{1. Năng lực tính toán số đặc trưng} & Nhầm lẫn công thức; xác định sai nhóm chứa trung vị, mốt; tính sai kết quả. & Tính đúng $\overline{x}$, xác định đúng nhóm chứa $M_e, M_o$; ráp đúng công thức. & Tính toán chính xác toàn bộ các số đặc trưng $\overline{x}, M_e, M_o, Q_1, Q_3, \Delta_Q$. & Phân tích sâu sắc, so sánh đa chiều và rút ra kết luận thực tiễn chính xác từ số liệu. \\
\hline
\textbf{2. Ứng dụng số (Mã 4.1NC1a)} & Không quan tâm đến việc bảo vệ dữ liệu cá nhân khi làm việc với số liệu. & Nhận biết được yêu cầu cần bảo mật thông tin cá nhân trong khảo sát số. & Thực hiện nghiêm túc việc ẩn danh hóa dữ liệu cá nhân trong báo cáo thống kê. & Tuyên truyền và hướng dẫn bạn học các quy tắc bảo mật quyền riêng tư trên môi trường số. \\
\hline
\textbf{3. Năng lực AI (Mã 11.A3.1)} & Chưa hiểu các khái niệm về đạo đức dữ liệu và quyền người dùng. & Nêu được quyền cơ bản của người dùng đối với dữ liệu của mình trong hệ thống AI. & Giải thích được nguy cơ định kiến thuật toán AI khi dữ liệu thống kê bị lệch. & Đánh giá và phản biện được các vấn đề đạo đức dữ liệu trong các ứng dụng AI thực tế. \\
\hline
\end{tabularx}
\end{center}

\end{document}
